\documentclass[10pt,a4paper]{amsart}
\usepackage[margin=19mm]{geometry}
\usepackage{amsmath,amssymb,mathtools}
\usepackage[scr=rsfs,cal=euler]{mathalfa}
\usepackage{xfrac}
\usepackage[unicode,hidelinks,bookmarks=false]{hyperref}
\newcommand{\ie}{i.e.\ }
\newcommand{\cf}{cf.\ }
\newcommand{\resp}{respectively\ }
\newcommand{\Subsection}{\subsection}
\providecommand{\nobreakdash}{\nobreak\hbox{-}}
\setlength{\emergencystretch}{2em}
\multlinegap0pt
\usepackage{amssymb}
\usepackage{float}
\newtheorem{proposition}{Proposition}
\title{An answer regarding automorphisms of finite abelian groups}
\author{Ryan McCulloch}
\date{Source: arXiv:2603.29299v2 (CC BY 4.0)}
\begin{document}
\maketitle

\noindent\emph{Source and licence.} An answer regarding automorphisms of finite abelian groups. Version arXiv:2603.29299v2 (CC BY 4.0), distributed by arXiv under the Creative Commons Attribution 4.0 International licence, http://creativecommons.org/licenses/by/4.0/, as recorded in the arXiv licence metadata for this exact version and shown on the abstract page https://arxiv.org/abs/2603.29299v2. Full licence notice: \texttt{licenses/arxiv-2603.29299-CC-BY-4.0.txt}.\par\smallskip

\noindent\textit{Editorial scope.} Counted unit: the article's proposition class (source0.tex lines 156--165). The article's complete original proof of it is delivered with it (source0.tex lines 167--171, one \texttt{proof} environment of 5 lines). No mathematics is written, repaired or re-derived: the statement and the proof are the article's own lines, in the article's own notation, and no retained line is substituted or re-flowed.  Retained uncounted as the source-local context the proof needs: lines 134--136; lines 152--154.  Nothing was omitted from any retained span, so the package carries no omission record.  No retained line contains a citation, so the wrapper carries no bibliography and no bibliography entry.  No retained line contains a figure environment, an image-inclusion command or drawing code, so no image is delivered and the package declares no asset dependency.  Editorial formatting only, in addition: an AMS \texttt{amsart} preamble that restates the article's own declarations and macros used by the retained lines, namely the environments \texttt{proposition} on separate counters, together with the standard packages those lines need.  The retained environments print in this document as Proposition 1 in order of appearance, whereas the article prints the same items with the numbering of its own full document.  The complete native source of the article as distributed in the arXiv e-print for this exact version is bundled unchanged as \texttt{sources/197506/source0.tex}.
\begin{proposition}\label{prop: form}
Let $G \cong \mathbb{Z}_{p^{e_1}} \times \cdots \times \mathbb{Z}_{p^{e_n}}$ be a finite abelian $p$-group, where $1 \leq e_1 \leq \dots \leq e_n$.  Then $$|\mathrm{Aut}(G)| = \prod_{k=1}^{n} (p^{d_k} - p^{k-1}) \prod_{j=1}^{n} (p^{e_j})^{n-d_j} \prod_{i=1}^{n} (p^{e_i-1})^{n-c_i+1},$$ where $d_r = \max\{s: e_s = e_r\}$ and $c_r = \min\{s: e_s = e_r\}$ for $r = 1, \dots, n$.
\end{proposition}

\begin{proposition}\label{prop: dc}
Let $G \cong (\mathbb{Z}_{p^{e_1}})^{k_1} \times \cdots \times (\mathbb{Z}_{p^{e_m}})^{k_m}$ be a finite abelian $p$-group, where $1 \leq e_1 < \dots < e_m$ and each $k_i \geq 1$.  Let $n = k_1 + \cdots + k_m$ and let $$d = \sum_{j=1}^{m-1} \Big( e_j k_j \big(\sum_{i=j+1}^m k_i \big) \Big)$$ and let $$c = \sum_{j=1}^{m} \Big( (e_j-1)k_j\big(\sum_{i=j}^m k_i \big) \Big).$$  Then $p^{((n-1)n/2) + d + c}$ is the largest $p$-power dividing $|\mathrm{Aut}(G)|$.
\end{proposition}

\begin{proposition}\label{prop: class}
Let $G$ be a finite abelian $p$-group.  
\begin{enumerate}
\item If $G$ is cyclic, then $|\mathrm{Aut}(G)|/|G| = (p-1)/p$.
\item If $G \cong \mathbb{Z}_{p} \times \mathbb{Z}_{p}$, then $|\mathrm{Aut}(G)|/|G| = (p-1)^2(p+1)/p$.
\item If $G \cong \mathbb{Z}_{p} \times \mathbb{Z}_{p^i}$, $i > 1$, then $|\mathrm{Aut}(G)|/|G| = (p-1)^2$.
\item If $G \cong \mathbb{Z}_{p} \times \mathbb{Z}_{p} \times \mathbb{Z}_{p}$, then $$|\mathrm{Aut}(G)|/|G| =(p-1)^3(p+1)(p^2+p+1).$$
\end{enumerate}
In all other cases, we have $|\mathrm{Aut}(G)|/|G|$ is an integer and is divisible by $p(p-1)^2$.
\end{proposition}

\begin{proof}
Items (1)--(4) are readily verified using Proposition \ref{prop: form}.  If $G$ is not cyclic, then by Proposition \ref{prop: form} we have $(p-1)^2$ divides $|\mathrm{Aut}(G)|$.  We work to show that in all cases other than items (1)--(4), we have $|\mathrm{Aut}(G)|/|G|$ is an integer and is divisible by $p$.

Use the same notation as in Proposition \ref{prop: dc} so that $|G| = p^a$ where $a = {\sum_{i=1}^m e_i k_i}$ and $n = {\sum_{i=1}^m k_i}$.  Note that $c \geq a-n$.  Also note that $(n-1)n/2 \geq n$ when $n \geq 3$ and $(n-1)n/2 = n-1$ when $n=1$ or $n=2$.  Suppose $n \geq 3$.  By Proposition \ref{prop: dc}, $|\mathrm{Aut}(G)|/|G|$ is an integer, and in this case (when $n \geq 3$), since $c + n \geq a$, we have $|\mathrm{Aut}(G)|/|G|$ coprime with $p$ implies $d=0$ and $c = a - n$.  Now $d=0$ implies $m=1$ and hence $k_1 = n \geq 3$.  Thus $c = a - n$ implies $e_1 = 1$, and we are in the situation of item (4).  Thus we can suppose $n \leq 2$.  Note that $n=1$ implies that $G$ is a cyclic, which is item (1), and so suppose $n=2$.  Hence $G \cong \mathbb{Z}_{p^{e_1}} \times \mathbb{Z}_{p^{e_2}}$ or $G \cong (\mathbb{Z}_{p^{e_1}})^2$.  In cases other than item (2) or (3), we have $e_1 > 1$.  If $G \cong \mathbb{Z}_{p^{e_1}} \times \mathbb{Z}_{p^{e_2}}$, then $m = 2$, and it follows that $d = e_1 \geq 2$.  Hence $c + d + n \geq a+2$ which implies $c + d + (n-1)n/2 \geq a+1$, and we have $|\mathrm{Aut}(G)|/|G|$ is an integer and is divisible by $p$.  If $G \cong (\mathbb{Z}_{p^{e_1}})^2$, then $m=1$ and $k_1=2$, and $c = 4(e_1-1)$ and $a-n = 2(e_1-1)$.  So $c \geq a-n + 2$, and hence $c + (n-1)n/2 \geq a+1$, and thus $|\mathrm{Aut}(G)|/|G|$ is an integer and is divisible by $p$. 
\end{proof}

\end{document}
